\section{Discussion}
\label{sec:discussion}

\paragraph{What the residual is.}
The target-inherent transfer residual (Section~\ref{sec:discovery}) is the
component of a guide's behavior that a source-trained model gets wrong in
the \emph{same way} across three nuclease screens whose own outcomes agree
only at $\rho \approx 0.73$. Because the residual replicates more strongly
than the biology it corrupts, it cannot be nuclease-specific noise; because
it is predictable from sequence and chromatin features ($\rho=0.78$
cross-screen, $0.86$ within-screen), it is not idiosyncratic measurement
error either. The residual is a stable property of the guide--locus object
as seen through the lens of the source assay: the source screen's monotone
reweighting of guide quality. Mechanistic candidates include local DNA
shape, repair-pathway context of the target locus, and sequence
composition effects on uptake or transcription in the source assay's
delivery system. The feature decomposition rules out any single simple
driver (max $|\rho|=0.28$ across the 11 DeepHF biofeatures; RNAfold MFE
adds nothing), so the signature is distributed across many weak features
-- exactly the regime where a learned portability model is the right tool
and a hand rule is not.

\paragraph{What practitioners should do differently.}
Three recommendations follow. First, \emph{never report} an efficacy
model's within-screen performance as evidence of portability: our gap-1
matrix shows within-family transfer ($0.43$--$0.50$) and cross-family
collapse ($0.03$) coexist for the same model. Second, when a target
library shares guides with any large public screen, a one-time
portability calibration (a few minutes of ridge regression against the
library's own residual) can recover most of the lost ranking signal on a
new assay (Table~9) -- the destination assay's labels are not needed.
Third, off-target scores should ship with calibration curves, but on
negative-scarce corpora the recalibration itself is starved: with
${\sim}$750 validation pairs, isotonic raises mean ECE ($0.326 \to
0.389$) and helps only one of five splits, and Platt is no better --
small validation slices cannot fit even a monotone map reliably. We
report the uncalibrated scores with ECE and let users calibrate on their
own larger holdouts.

\paragraph{Why the pegRNA transfer negative matters.}
The naive hypothesis that Cas9 efficacy motifs transfer to pegRNA design
is attractive -- the models share a substrate -- and it fails at every
sample size we tested (Figure~3). The honest reading is that prime-editing
efficiency is dominated by the repair outcome of the flap intermediate,
not by the cleavage-nick kinetics that Cas9 efficacy models capture. This
argues for investment in pegRNA-specific screens (PRIDICT-scale and
larger) rather than transfer shortcuts, and it makes the pegRNA gap a
data problem, not a modeling problem.

\paragraph{Relation to domain adaptation theory.}
Proposition~\ref{prop:indisc} gives the classical sufficient condition
(affine-equivalent conditional laws) under which transfer is free, and
our gap-1 measurements show how violently it fails between assay
families. The discovery adds structure to that failure: the error is not
an unstructured shift but a stable, learnable, guide-level object. In
domain-adaptation terms, the residual behaves like a shared labeling
function artifact with low cross-screen variability -- the exact case
where a cheap correction beats retraining. We are not aware of prior work
identifying or exploiting this object in genomics.

\paragraph{Scope of claims.}
Every strong claim in this paper is confined to its measured regime:
PORTA is validated between Doench-family source models and the three
DeepHF nuclease screens; the calibration results hold on crisprSQL's
17-study corpus with its 94\% positive prevalence; the pegRNA negative is
at PRIDICT2 scale with motif-level features. Extrapolation beyond these
regimes (base editors, Cas12a, non-U2OS chromatin) is future work, not
implied.

\paragraph{Negative-results ledger.}
This project produced more informative negatives than positives, and we
consolidate them here because each is load-bearing for the next
practitioner.
(i)~\emph{Pooling datasets does not fix cross-dataset collapse} (gap 1):
pooling Doench FC+RES with DeepHF degrades transfer in both directions.
(ii)~\emph{Cas9$\to$pegRNA transfer never beats scratch} (gap 3):
not at $n=250$, $1000$, or $4000$, in HEK or K562.
(iii)~\emph{Three independent thermodynamic channels add nothing} to
pegRNA numerics (gap 3): ViennaRNA MFE ($+0.0002$), seqfold ensemble
$\Delta G$ and primer3 $T_m$ ($\le 10^{-4}$ Spearman).
(iv)~\emph{Isotonic recalibration hurts} on validation slices of this
size (gap 2): mean ECE $0.326 \to 0.389$, helping 1 of 5 splits.
(v)~\emph{PAM and chromatin channels do not robustly beat seq-only}
(gaps 4--5): five-seed replication places every condition within noise
($0.590 \pm 0.094$ seq-only; PAM-alone worst at $0.528 \pm 0.119$).
(vi)~\emph{Study routing is not estimable} on crisprSQL (gap 2): a
leave-one-out router collapses to always-MIT, the prevalence rule
underperforms it, and an oracle gains $0.001$; apparent per-study wins
rest on as few as two positives.
(vii)~\emph{Initialization is a first-class variance source} (gap 2):
on a byte-identical split, torch init alone spans AUROC
$0.530$--$0.606$ and ECE $0.209$--$0.389$.
We report these because each was, at some point, our own leading
hypothesis, and each fell to a control we initially considered
pedantic. The pattern generalizes: on negative-scarce,
study-confounded corpora, the honest effect size of most modeling
choices is smaller than the evaluation noise.
